\section{Detailed Statistics and Per-Trial Data}
\label{appendix:per_trial}

\noindent\textit{This appendix provides per-stratum breakdowns and detailed statistics
for individual controllers. The full 13-controller benchmark results
(Table~II equivalent) are given in Table~\ref{tab:full_results} below.}

\subsection{MPPI: Per-Stratum Breakdown (50 Trials)}

\begin{table}[t]
\centering
\caption{MPPI hardware results by stratum.}
\begin{tabular}{lccccc}
\toprule
\textbf{Difficulty} & \textbf{Trials} & \textbf{SR} & \textbf{Mean (s)} & \textbf{Med (s)} & \textbf{Max (s)} \\
\midrule
Center & 10 & 100\% & 1.82 & 1.29 & 4.59 \\
Mid-rail & 16 & 100\% & 1.99 & 1.59 & 4.59 \\
Near-wall & 14 & 100\% & 2.33 & 2.19 & 7.04 \\
At-wall-opposite & 5 & 100\% & 2.52 & 2.69 & 4.19 \\
At-wall-same & 5 & 100\% & 1.56 & 1.19 & 2.34 \\
\midrule
\textbf{Overall} & \textbf{50} & \textbf{100\%} & \textbf{2.06} & \textbf{1.67} & \textbf{7.04} \\
\bottomrule
\end{tabular}
\end{table}

MPPI's median settling varies little across strata: it barely
moves with the starting position, from 1.19\,s at at-wall-same and 1.29\,s at
center to 2.69\,s at at-wall-opposite, a spread of only about 1.5\,s despite
the very different maneuvering room.

\subsection{Complete Benchmark Results (All 13 Controllers)}

Table~\ref{tab:full_results} reports the definitive 50-trial hardware
benchmark for all 13 deployed controllers, sorted by hardware success
rate then by mean settling time (the same controllers as Table~II of the main
paper, extended here with the bootstrap settling-time confidence band and
mean per-step compute time). SR 95\% confidence intervals are Wilson
score intervals (matching Appendix~\ref{appendix:protocol} and the main paper);
settling 95\% confidence intervals are bootstrap over the successful
trials only, and the per-step compute time excludes sensor I/O. Two controllers (MPPI, PPO~(WM)) reach
100\% SR; PPO~(BZ+DR) and PD+WO reach 96\%; the rest span
80--92\%. The classical feedback laws, MPC+WO, and the on-policy RL policies PPO~(BZ+DR) and TRPO~(DR) run in under 1\,ms; IQL, NMPC, PPO-WM, and the off-policy RL policies
(TD3, TQC, SAC) sit in the 2--8\,ms range; and only MPPI ($\sim$40\,ms)
approaches the 50\,ms deadline.

\begin{table*}[t]
\centering
\caption{Complete hardware benchmark: 50 trials per controller, 20\,Hz,
stratified initial conditions. SR intervals are 95\% Wilson score intervals; settling intervals are 95\% bootstrap.}
\label{tab:full_results}
\begin{tabular}{llcccccc}
\toprule
\textbf{Controller} & \textbf{Paradigm} & \textbf{SR (\%)} & \textbf{SR 95\% CI}
  & \textbf{Settle (s)} & \textbf{Settle 95\% CI} & \textbf{RMS Effort}
  & \textbf{Compute (ms)} \\
\midrule
MPPI & Real-data & 100 & [93, 100] & 2.06 & [1.71, 2.44] & 0.354 & 40.0 \\
PPO (WM) & Real-data & 100 & [93, 100] & 6.84 & [5.07, 8.95] & 0.492 & 6.7 \\
PPO (BZ+DR) & Sim-trained RL & 96 & [87, 99] & 3.81 & [2.93, 4.78] & 0.498 & 0.59 \\
PD+WO & Classical & 96 & [87, 99] & 5.80 & [4.08, 7.76] & 0.580 & 0.05 \\
SMC+WO & Classical & 92 & [81, 97] & 8.93 & [7.33, 10.64] & 0.562 & 0.06 \\
LQR+WO & Classical & 92 & [81, 97] & 10.21 & [7.76, 12.79] & 0.625 & 0.33 \\
IQL & Real-data & 92 & [81, 97] & 10.62 & [8.54, 12.76] & 0.395 & 2.0 \\
TD3 (DR) & Sim-trained RL & 92 & [81, 97] & 11.45 & [9.16, 13.87] & 0.680 & 7.9 \\
NMPC & Predictive & 90 & [79, 96] & 7.48 & [5.19, 10.08] & 0.608 & 5.7 \\
TRPO (DR) & Sim-trained RL & 90 & [79, 96] & 8.09 & [6.01, 10.43] & 0.508 & 0.67 \\
TQC (DR) & Sim-trained RL & 86 & [74, 93] & 9.77 & [7.25, 12.39] & 0.663 & 7.9 \\
SAC (DR) & Sim-trained RL & 82 & [69, 90] & 10.10 & [7.80, 12.42] & 0.588 & 7.8 \\
MPC+WO & Predictive & 80 & [67, 89] & 10.81 & [8.04, 13.55] & 0.623 & 0.47 \\
\bottomrule
\end{tabular}
\end{table*}

\begin{figure}[t]
  \centering
  \includegraphics[width=0.95\columnwidth]{figures/hw_sr_ranking.eps}
  \caption{Hardware success rate ranking across all controllers
    (50 trials each, 20~Hz, stratified ICs). Bars are colored by paradigm
    (real-data, classical, sim-trained RL, and predictive), matching
    Table~\ref{tab:full_results}. MPPI and PPO~(WM)
    achieve 100\% SR; PPO~(BZ+DR) and PD+WO reach 96\%;
    SMC+WO, LQR+WO, IQL~(offline), and TD3~(DR) reach 92\%.}
  \label{fig:hw_ranking}
\end{figure}

\subsection{Failed Controller Variants (Excluded from Main Results)}

Table~\ref{tab:failed_variants} lists controller variants excluded from
the main benchmark. The first six did not reach our 50\% SR deployment
threshold on hardware; PPO (base) and PPO (DR) are the un-augmented and
DR-only PPO ablations (the deployed form is PPO~(BZ+DR)). They are reported here for completeness
because their failure modes are the most instructive design lessons of the
study: in every case the hardware collapse traces to a specific
missing capability: absence of boundary handling (PD, LQR, SMC),
a constraint formulation incompatible with the velocity-commanded actuator
(MPC QP), or the sim-to-real gap left by not modeling the sensor blind zone
or by omitting domain randomization (PPO base). Adding the wall-override
heuristic or the blind-zone+DR training configuration lifts each of these to
the corresponding +WO / +BZ+DR row of Table~\ref{tab:full_results}, which is
why the deployed benchmark uses those augmented variants.

\begin{table*}[t]
\centering
\caption{Excluded and ablation controller variants; the first six fall
         below the 50\% deployment threshold on hardware (the two PPO
         ablations reach 60\% and 90\%). Included for completeness,
         these demonstrate the critical importance of constraint handling.
         Sim SR is reported for the first-order-lag deployment regime; the
         S-curve development regime scores these variants higher
         (Appendix~\ref{appendix:rl_training}).}
\label{tab:failed_variants}
\begin{tabular}{lccp{7cm}}
\toprule
\textbf{Controller} & \textbf{HW SR} & \textbf{Sim SR} & \textbf{Failure Cause} \\
\midrule
PD (vanilla) & 18\% & 84\% & Positive-feedback runaway: ball drifts, controller saturates, cart hits wall \\
LQR (basic) & 20\% & 84\% & No boundary awareness; nominal feedback runs out of rail at the limits \\
SMC (basic) & 12\% & 14\% & Wide boundary layer ($\phi=0.85$) + no WO; controller too passive at limits \\
MPC (QP constraints) & 44\% & 34\% & QP infeasibility near walls; constraint formulation incompatible with velocity interface \\
LQR (Q-center) & 28\% & 100\% & Q-weighting on cart position ($K_x=3.31$) structurally insufficient for constraint handling \\
SMC (full) & 24\% & 100\% & Centering + barrier terms do not overcome sensor-offset and blind-zone effects on hardware \\
PPO (base, no DR) & 60\% & 100\% & Sim-to-real gap: blind zone not modeled; DR-induced robustness absent \\
PPO (DR, no BZ) & 90\% & 100\% & Blind-zone (VL53L0X saturation) not modeled; residual sim-to-real gap (deployed form adds BZ) \\
\bottomrule
\end{tabular}
\end{table*}

\subsection{Statistical Significance}

We run pairwise Mann-Whitney $U$ rank tests with Bonferroni correction
($\alpha = 0.05/78$ for 13 controllers) on the settling times of
successful trials. Eighteen of the 78 pairs are significant; because these
pairwise tests share trial data (many pairs involve a common controller),
the Bonferroni correction is conservative here, so this count likely
understates the number of true differences. MPPI is
significantly faster than the slower nine controllers (PPO~(WM), LQR+WO, SMC+WO,
TD3, IQL, TRPO, TQC, SAC, MPC+WO) but not than the other fast, reliable
controllers (PPO~(BZ+DR), PD+WO, NMPC), which it beats by a large
but sub-threshold margin (Cohen's $d$ of 0.70 to 0.90; e.g.\ 0.90 versus
NMPC, 0.70 versus PPO~(BZ+DR)). Those three, together with PPO~(WM), form a
fast cluster with no significant internal differences; the mid- and
low-table controllers likewise separate into loose clusters with few
significant differences inside each. Under Fisher's exact test with the
same correction, no pairwise success-rate difference is significant, so
the reliability ordering is practical rather than statistically certified.

\begin{figure}[t]
  \centering
  \includegraphics[width=0.95\columnwidth]{figures/exp1_significance_heatmap.eps}
  \caption{Pairwise settling-time significance for the 13 benchmarked
    controllers (Mann-Whitney $U$ rank test on successful trials, Bonferroni
    correction $\alpha = 0.05/78$; 18 of 78 pairs significant). Green cells
    mark significant differences.}
  \label{fig:significance}
\end{figure}

\subsection{Initial Condition Protocol}

We displace the ball from center with a brief cart ``jolt/shake'', noting
its direction post-hoc from the first sensor reading; the cart is moved
automatically by the evaluation script, with no manual ball placement during
trials. For the two at-wall strata the ball side is targeted, and the realized
cart position and side (read per trial from the recorded initial state) vary
somewhat from the nominal targets, as detailed in
Appendix~\ref{appendix:protocol}.
The five-stratum, 50-trial schedule
(Center, Mid-rail, Near-wall, and the two at-wall strata) is specified in
Appendix~\ref{appendix:protocol}.

\subsection{Control Effort and Trajectory Shape}

The qualitative differences behind these numbers show up in three views. The
control-effort distributions (Fig.~\ref{fig:effort_boxplot}) put MPPI and IQL
at the smooth, low-effort end and TD3 and TQC at the aggressive, high-effort
end, matching their RMS-effort figures in Table~\ref{tab:full_results}.
A per-controller gallery of representative ball-angle recoveries
(Fig.~\ref{fig:multi_traj}, one panel per controller, each showing that
controller's median-settling trial) shows the same contrast in recovery
shape: MPPI and PD+WO settle within a few seconds with little overshoot,
while the other wall-override classical controllers (LQR+WO, SMC+WO, MPC+WO)
and the slower simulation-trained agents (SAC, TQC, TD3) take longer and hunt
more before settling. Because each
panel uses a different representative trial, these are illustrations of
typical recovery shape, not a matched head-to-head comparison. The command profiles (Fig.~\ref{fig:action_profiles}) trace one clean recovery
per controller; once the ball is settled the command drops to a small residual
(last-second action RMS around 0.05--0.23 for these representative trials, well
below the $\pm1$ saturation), so these clean-recovery exemplars settle to a
small residual command.

\begin{figure}[t]
  \centering
  \includegraphics[width=0.95\columnwidth]{figures/effort_boxplot.eps}
  \caption{RMS control effort distribution over all trials.
    MPPI and IQL produce the smoothest, lowest-effort control signals.
    TD3 and TQC show the highest effort, indicating aggressive,
    oscillatory control.}
  \label{fig:effort_boxplot}
\end{figure}

\begin{figure}[t]
  \centering
  \includegraphics[width=0.95\columnwidth]{figures/compute_time_bars.eps}
  \caption{Mean per-step compute time by controller. MPPI ($\sim$40\,ms)
    runs close to the 50\,ms deadline (red dashed line). IQL, the off-policy
    RL policies (TD3, TQC, SAC), PPO-WM, and NMPC form a second tier at
    2--8\,ms; the classical controllers, MPC, and the on-policy RL policies
    execute in under 1\,ms.}
  \label{fig:compute_time}
\end{figure}

\begin{figure*}[t]
  \centering
  \includegraphics[width=0.95\textwidth]{figures/multi_controller_trajectory.eps}
  \caption{Representative ball-angle recovery for each of the 13 benchmark
    controllers. Each panel shows that controller's successful trial closest
    to its own median settling time (orange dotted line: settling time;
    green band: the $\pm$0.01\,rad settling window). Because each panel uses
    that controller's own representative trial, the initial conditions differ
    across panels, so this is a per-controller view of typical recovery shape
    rather than a matched head-to-head. MPPI and PD+WO recover within a few
    seconds; LQR+WO, SMC+WO, MPC+WO and the slower simulation-trained
    agents (SAC, TQC, TD3) settle later.}
  \label{fig:multi_traj}
\end{figure*}

\begin{figure*}[t]
  \centering
  \includegraphics[width=0.95\textwidth]{figures/action_profiles.eps}
  \caption{Control action profiles for all 13 benchmark controllers.
    For each controller the shown trial is selected by a single uniform
    rule: among successful trials that settled reasonably early
    (settling time $\le 1.5\times$ that controller's median) and have at
    least a 1\,s settled tail (episode length $\ge$ settling $+ 1$\,s),
    keep those whose final-1\,s action RMS is within 10\% of that
    controller's best, and pick the earliest-settling of those, i.e.\ a
    clean, early-settling exemplar rather than a barely-settled outlier.
    Shaded green band = the final 1\,s window; orange dotted line = the
    trial's settling time. Last-1\,s action RMS is 0.05--0.23 across
    controllers, i.e.\ commands are near zero once settled. Generated
    from \texttt{paper/data/exp1\_hardware/consolidated.json} by
    \texttt{regenerate\_extra\_figures.py}.}
  \label{fig:action_profiles}
\end{figure*}
